% TOP_PROBLEM: {"id": 2991, "title": "Kirby Problem 4.115", "category_id": 7, "category": {"id": 7, "name": "topology", "display_name": "Topology", "description": "Properties preserved under continuous deformations.", "slug": "topology", "order_index": 7, "created_at": "2026-07-31T15:26:25.671Z"}, "status": "open", "problem_number": "KP-4.115", "set_id": 11}
% FIRST_POSTED: 2026-10-01
% ENTRY_KIND: statement_only
% UnsolvedMath ID 2991; source catalogue code KP-4.115
% !TeX program = lualatex
% Definition and sources only; this file is not an LLM solution attempt.
\documentclass[11pt,a4paper]{article}
\usepackage[margin=25mm]{geometry}
\usepackage{fontspec}
\setmainfont{lmroman10-regular.otf}[BoldFont=lmroman10-bold.otf,
 ItalicFont=lmroman10-italic.otf,BoldItalicFont=lmroman10-bolditalic.otf]
\usepackage{amsmath,amssymb,mathtools}
\usepackage{unicode-math}
\setmathfont{latinmodern-math.otf}
\AtBeginDocument{\let\setminus\smallsetminus}
\usepackage{xurl,hyperref}
\hypersetup{colorlinks=true,urlcolor=blue}
\setcounter{secnumdepth}{0}
\setlength{\parindent}{0pt}
\setlength{\parskip}{5pt}
\setlength{\emergencystretch}{3em}
\begin{document}
\section*{Kirby Problem 4.115}\label{2991}
{\small UnsolvedMath ID 2991 · KP-4.115 · Topology · Kirby's Problems in Low-Dimensional Topology}
\subsection{Definitions and mathematical statement}
A $(g;k_1,k_2,k_3)$ trisection of a closed oriented smooth four-manifold $X$ is a decomposition into four-dimensional one-handlebodies $X_i\cong\natural^{k_i}(S^1\times B^3)$ whose pairwise intersections are three-dimensional genus-$g$ handlebodies and whose triple intersection is a closed genus-$g$ surface. It is balanced when $k_1=k_2=k_3$.

Two trisections of the same $X$ are diffeomorphic if an orientation-preserving self-diffeomorphism carries their sectors to one another. They are isotopic if such a carrying diffeomorphism is the endpoint of a smooth ambient isotopy starting at the identity. Allow a permutation of sector labels consistently in both notions; a mere relabeling is not an example of inequivalence.

(a) Find a closed smooth four-manifold with two diffeomorphic trisections that are not isotopic. (b) Find a closed simply connected smooth four-manifold with two balanced trisections of the same genus that are not diffeomorphic.

In (a), applying a diffeomorphism not isotopic to the identity to a trisection is only a candidate construction: some different diffeomorphism isotopic to the identity could carry the initial trisection to its image. All such possible carrying maps must be excluded. In (b), even maps in other components of the diffeomorphism group must be excluded. The shared genus prevents distinguishing the decompositions only by their central-surface genus, while simple connectivity rules out explanations based solely on different presentations of a nontrivial fundamental group.
\subsection{Short English statement}
Find trisections distinguished by isotopy despite diffeomorphism, or balanced equal-genus trisections distinguished even by diffeomorphism.
\subsection{Sources}
\begin{itemize}
\item[S1] ULAM.AI and the UnsolvedMath Contributors, supplied problems.json, record 2991 (KP-4.115), Kirby's Problems in Low-Dimensional Topology. Catalogue identity and source transcription; CC BY 4.0.
\url{https://huggingface.co/datasets/ulamai/UnsolvedMath}.
\item[S2] K3: A New Problem List in Low-Dimensional Topology, Problem 4.115. Expanded formulation of the supplied catalogue statement.
\url{https://math.berkeley.edu/sites/default/files/surv-295-ruberman-watermarked-author-pdf.pdf}.
\end{itemize}
\end{document}
